\documentclass[12pt]{article}
\usepackage{cwilliams-standard}

\setclass{\NLP}
\settitle{Homework 1}

\definecolor{tableheader}{rgb}{0.18, 0.32, 0.55}
\definecolor{rowalt}{rgb}{0.93, 0.95, 0.98}

\begin{document}

\maketitlepage

In this report, I will go over my steps in tackling the first homework on NLP. These tasks included: 
Text-Embedding and Topic Modeling. It was also unstated (but clearly) necessary that we needed to do 
preprocessing on the corpora given.

For the first (zeroth, $0^{th}$) task, I did preprocessing. For each subsequent task, I attempted to do what I 
believed was appropriate preprocessing for said task (as we talked at length in class about how 
preprocessing was approached differently depending on what the problem you are dealing with is). I will 
do my best to explain each of the methods of preprocessing I had considered, and which ones I decided on 
using for each task and why.

\begin{aside}{Methodology}
For the two main tasks, I will explain my methodology as we go. 
\end{aside}

\begin{important}{Dataset}
    For the \textit{entire} report/assignment, I did not split the dataset up into its constituent parts 
    (train, validation, test) but, instead, \textbf{combined} them (as the assignment said to do, Pg. 1). Then, after combination 
    they are used in 5-fold cross validation. Stated for clarity.
\end{important}

\begin{aside}{Code}
    The code will be submitted on Canvas with this report PDF. If, for some reason, the code does not 
    behave or load properly, the file is uploaded to Google Drive and there will be a link to the code 
    in the Code section of the report.
\end{aside}


\newpage
\section*{Preprocessing}

Firstly, for preprocessing I laid out these options of preprocessing the data to give myself as many options 
of cleaning the data as possible.

\begin{table}[H]
    \centering
    \definecolor{tableheader}{rgb}{0.18, 0.32, 0.55}
    \definecolor{rowalt}{rgb}{0.93, 0.95, 0.98}
    \renewcommand{\arraystretch}{1.4}
    \begin{tabular}{>{\bfseries\centering\arraybackslash}p{0.6cm}
                    p{5.2cm}
                    p{7.8cm}}
        \toprule
        \rowcolor{tableheader}
        \color{white}\# & \color{white}\textbf{Method} & \color{white}\textbf{Description} \\
        \midrule
        \rowcolor{rowalt} 1 & Expand Contractions      & \texttt{"can't"} $\to$ \texttt{"cannot"} \\
                          2 & Lowercase                & Normalize all text to lowercase \\
        \rowcolor{rowalt} 3 & Remove URLs \& Mentions  & Strip \texttt{http://...} and \texttt{@username} tokens \\
                          4 & Strip Hashtag Symbol     & \texttt{\#word} $\to$ \texttt{word} (keep the term) \\
        \rowcolor{rowalt} 5 & Remove Emoji             & Drop all emoji characters \\
                          6 & Remove Punctuation       & Strip non-alphanumeric characters \\
        \rowcolor{rowalt} 7 & Remove Digit-only Tokens & Drop tokens consisting entirely of digits \\
                          8 & Remove Stop Words        & Filter standard English stop words \\
        \rowcolor{rowalt} 9 & POS-aware Lemmatization  & WordNet lemmatization guided by POS tags \\
        \bottomrule
    \end{tabular}
    \caption{Preprocessing Options}
    \label{tab:preprocessing}
\end{table}

These options were taken from class with one exception: \# 9, POS-aware Lemmatization. 
This was suggested by Claude Code (within VSCode) to tag parts of speech in the preprocesing stage. 
This specific lemmatization is aware of parts of speech (i.e., understands that "running" 
is a verb and will convert that to its root, "run") but does not pass along any additional 
information about the word to the frequency-based embedder.

Each preprocessing option that was used in each task will be stated.

\phantomsection\section*{Text-Embedding}

For each subtask under this section, each model (Random Forest, SVM, Decision Tree) was tested
on each embedding schema. Here are the hyperparameters I used for each model.

\begin{table}[H]
    \centering
    \definecolor{tableheader}{rgb}{0.18, 0.32, 0.55}
    \definecolor{rowalt}{rgb}{0.93, 0.95, 0.98}
    \renewcommand{\arraystretch}{1.4}
    \begin{tabular}{p{5.0cm} p{8.0cm}}
        \toprule
        \rowcolor{tableheader}
        \color{white}\textbf{Hyperparameter} & \color{white}\textbf{Value} \\
        \midrule
        \multicolumn{2}{l}{\cellcolor{rowalt!60}\textbf{Random Forest}} \\
        \rowcolor{rowalt} \texttt{n\_estimators}  & 200 \\
                          \texttt{random\_state}  & 42 \\
        \rowcolor{rowalt} \texttt{class\_weight}  & \texttt{"balanced"} \\
                          \texttt{n\_jobs}        & $-1$ (use all cores) \\
        \midrule
        \multicolumn{2}{l}{\cellcolor{rowalt!60}\textbf{SVM}} \\
        \rowcolor{rowalt} \texttt{kernel}         & \texttt{"rbf"} \\
                          \texttt{C}              & 1.0 \\
        \rowcolor{rowalt} \texttt{gamma}          & \texttt{"scale"} \\
                          \texttt{random\_state}  & 42 \\
        \rowcolor{rowalt} \texttt{class\_weight}  & \texttt{"balanced"} \\
        \midrule
        \multicolumn{2}{l}{\cellcolor{rowalt!60}\textbf{Decision Tree}} \\
        \rowcolor{rowalt} \texttt{criterion}      & \texttt{"gini"} \\
                          \texttt{random\_state}  & 42 \\
        \rowcolor{rowalt} \texttt{class\_weight}  & \texttt{"balanced"} \\
        \bottomrule
    \end{tabular}
    \caption{Classifier hyperparameters used across all embedding tasks.}
    \label{tab:hyperparameters}
\end{table}

 As stated in \textit{Hitchhikers Guide to the Galaxy} by Douglas Adams, 42 is the answer to everything.

\section{Frequency-based Embeddings} \label{freq_embed}

For this portion of the assignment, I was curious which frequency embedder would perform better. So, 
I used \textit{both} Bag of Words and TF-IDF to compare the two methods and see which resulted in 
better performance. Both embedders received the works of preprocessing (all available methods described above) 
before being fed the data.

Bag of Words represents each document in the corpora as a vector of word counts across the vocabulary. 
Words that appear more often in a document get higher values, but all words are treated equally regardless 
of how common they are across the entire corpus. This does allow uninformative filler words to have a higher 
weight (not stop words, as those were already removed in preprocessing). So, an update to that method is 
TF-IDF. This method weighs each word by how ``distinctive'' it is. Within document frequency is weighted 
heavy, but if a specific word appears often across the entire corpus, then it is reduced in importance. 
This method is looking for informative words within specific documents rather than the entire corpus.

I believed that by comparing the two methods I could see how much more effective TF-IDF was than Bag 
of Words for frequency embedding, potentially seeing a decent boost to model performance. However, the 
boost was not nearly as significant as I thought. 

Bag of Words only underperformed by a percentage at most, and actually beat TF-IDF in F1 score (by 0.0009)
in SVM. It also appears that Bag of Words beats TF-IDF in all models, within F1. 
This tiny difference surprised me and made me check my approach to make sure there wasn't a mistake, but 
it appears that frequency embedding is not as superior as I anticipated it to be.

A potential explanation of this performance is that stop word removal during preprocessing already eliminated 
a decent amount of globally common words that TF-IDF is designed to take into consideration, which may 
have reduced TF-IDF's advantage that I believed it would have.

\begin{figure}[H]
    \centering
    \begin{subfigure}[t]{0.48\textwidth}
        \centering
        \includegraphics[width=\textwidth]{images/results_bow.png}
        \caption{Bag of Words}
    \end{subfigure}
    \hfill
    \begin{subfigure}[t]{0.48\textwidth}
        \centering
        \includegraphics[width=\textwidth]{images/results_tfidf.png}
        \caption{TF-IDF}
    \end{subfigure}
    \caption{Frequency-based embedding results.}
    \label{fig:freq-embeddings}
\end{figure}

\begin{table}[H]
    \centering
    \definecolor{tableheader}{rgb}{0.18, 0.32, 0.55}
    \definecolor{rowalt}{rgb}{0.93, 0.95, 0.98}
    \renewcommand{\arraystretch}{1.4}
    \begin{tabular}{l
                    >{\centering\arraybackslash}p{1.8cm}
                    >{\centering\arraybackslash}p{1.8cm}
                    >{\centering\arraybackslash}p{1.8cm}
                    >{\centering\arraybackslash}p{1.8cm}
                    >{\centering\arraybackslash}p{1.8cm}
                    >{\centering\arraybackslash}p{1.8cm}}
        \toprule
        \rowcolor{tableheader}
        \color{white}\textbf{Classifier}
            & \color{white}\textbf{Acc.}
            & \color{white}\textbf{Acc. $\sigma$}
            & \color{white}\textbf{F1}
            & \color{white}\textbf{F1 $\sigma$}
            & \color{white}\textbf{Macro F1}
            & \color{white}\textbf{MF1 $\sigma$} \\
        \midrule
        \multicolumn{7}{l}{\cellcolor{rowalt!60}\textbf{Bag of Words}} \\
        \rowcolor{rowalt} Random Forest  & 0.6364 & 0.0143 & 0.5457 & 0.0163 & 0.6212 & 0.0139 \\
                          SVM            & \cellcolor{green!13}\textbf{0.6381} & 0.0118 & \cellcolor{green!13}\textbf{0.5619} & 0.0151 & \cellcolor{green!13}\textbf{0.6266} & 0.0109 \\
        \rowcolor{rowalt} Decision Tree  & 0.6359 & 0.0128 & 0.5456 & 0.0150 & 0.6208 & 0.0123 \\
        \midrule
        \multicolumn{7}{l}{\cellcolor{rowalt!60}\textbf{TF-IDF}} \\
        \rowcolor{rowalt} Random Forest  & 0.6270 & 0.0109 & 0.5052 & 0.0249 & 0.6029 & 0.0147 \\
                          SVM            & \cellcolor{green!13}\textbf{0.6438} & 0.0070 & \cellcolor{green!13}\textbf{0.5610} & 0.0106 & \cellcolor{green!13}\textbf{0.6305} & 0.0062 \\
        \rowcolor{rowalt} Decision Tree  & 0.6194 & 0.0080 & 0.4905 & 0.0233 & 0.5933 & 0.0122 \\
        \bottomrule
    \end{tabular}
    \caption{Frequency-based embedding results (5-fold CV). Acc.\ = Accuracy, MF1 = Macro F1, $\sigma$ = std.\ dev.}
    \label{tab:freq-results}
\end{table}

Shown in this table, SVM wins. SVM is meant for higher dimensional tasks, and text 
embedding in general falls into that category as there is not any one, clean decision boundary that 
"classifies" words. Also, the curse of dimensionality punishes tree-based methods like the Decision Tree 
and more primarily Random Forest \cite{SVMTextCat-Joachims}.
BOW and TF-IDF both produce high dimensionality, sparse vectors and SVMS are designed for higher dimension 
problems, finding a maximum-margin hyperplane with support vectors. \cite{redshift_svm_highdim} 
\cite{RFHighDim-Genuer}
However, ultimately, all the models performed very similarly meaning in general they were all gleaming 
the same amount of information from the preprocessed data. So, we should try another embedding technique!

\section{Word-based Embeddings} \label{word_embed}

For task 2, Word-based embeddings, I used 5 of the 8 available preprocessing methods: (2) Lowercase, 
(3) Remove URLs \& Mentions, (4) Strip Hashtag Symbol, (8) Remove Stop Words, and (9) POS-aware 
Lemmatization. The reasoning for this is that I used Word2Vec embedding, specifically the 
\textbf{English Wikipedia} (enwiki\_upos\_skipgram\_300\_2\_2021) \cite{fares-etal-2017-word}. 
They state on the website \cite{WebVectors-Fares2017} that the preprocesing used is as follows: 

\begin{center}
    \textit{split into sentences, tokenized, lemmatized, PoS-tagged, multi-word named entities extracted, 
    stop words removed}
\end{center}

To best fit to what the model was already trained on, I did the preprocessing that would get me closest 
to what the model already knew. I chose this particularly Word2Vec embedding due to how clear the preprocessing 
was and that it was the newest. I had emailed Dr. Garg about taking into consideration the 
SimLex999 and Google Analogy scores, but he had said to not take it into consideration for this 
assignment.

I chose Word2Vec over GloVe because I understood it better than GloVe (and it was also easier to implement).

Word2Vec obtains ``meanings'' of words based on vector representations of each word in relation to 
every other word. As the ``meaning'' of one word changes, every other meaning shifts as well. This representation 
of high dimensional vectors that contain meaning made the most sense to me and, given the high dimensionality 
of the vectors, believed that SVM would perform well once again. 

GloVe does the same process, but takes into account the position of the words within each document. Word2Vec 
does not consider the positioning of the words, so it simply loses out on that information. This would, in 
theory make GloVe the better word-based embedding approach.

After applying the preprocessing, I trained the 3 separate models on PCA and t-SNE (t-distributed 
stochastic neighbor embedding) dimensionality reduction and compared the 6 models. 

\begin{figure}[H]
    \centering
    \includegraphics[width=\textwidth]{images/results_word2vec.png}
    \caption{Word-based embedding results.}
    \label{fig:word2vec}
\end{figure}

\begin{table}[H]
    \centering
    \renewcommand{\arraystretch}{1.4}
    \begin{tabular}{l
                    >{\centering\arraybackslash}p{1.8cm}
                    >{\centering\arraybackslash}p{1.8cm}
                    >{\centering\arraybackslash}p{1.8cm}
                    >{\centering\arraybackslash}p{1.8cm}
                    >{\centering\arraybackslash}p{1.8cm}
                    >{\centering\arraybackslash}p{1.8cm}}
        \toprule
        \rowcolor{tableheader}
        \color{white}\textbf{Classifier}
            & \color{white}\textbf{Acc.}
            & \color{white}\textbf{Acc. $\sigma$}
            & \color{white}\textbf{F1}
            & \color{white}\textbf{F1 $\sigma$}
            & \color{white}\textbf{Macro F1}
            & \color{white}\textbf{MF1 $\sigma$} \\
        \midrule
        \multicolumn{7}{l}{\cellcolor{rowalt!60}\textbf{PCA}} \\
        \rowcolor{rowalt} Random Forest  & \cellcolor{green!13}\textbf{0.6036} & 0.0117 & 0.5323 & 0.0171 & 0.5941 & 0.0123 \\
                          SVM            & 0.6027 & 0.0156 & \cellcolor{green!13}\textbf{0.5656} & 0.0192 & \cellcolor{green!13}\textbf{0.5996} & 0.0157 \\
        \rowcolor{rowalt} Decision Tree  & 0.5392 & 0.0100 & 0.5228 & 0.0138 & 0.5386 & 0.0101 \\
        \midrule
        \multicolumn{7}{l}{\cellcolor{rowalt!60}\textbf{t-SNE}} \\
        \rowcolor{rowalt} Random Forest  & \cellcolor{green!13}\textbf{0.5588} & 0.0211 & 0.5318 & 0.0224 & \cellcolor{green!13}\textbf{0.5573} & 0.0211 \\
                          SVM            & 0.5384 & 0.0150 & 0.4853 & 0.0235 & 0.5326 & 0.0133 \\
        \rowcolor{rowalt} Decision Tree  & 0.5568 & 0.0197 & \cellcolor{green!13}\textbf{0.5450} & 0.0255 & 0.5564 & 0.0200 \\
        \bottomrule
    \end{tabular}
    \caption{Word2Vec embedding results by dimensionality reduction (5-fold CV). Acc.\ = Accuracy, MF1 = Macro F1, $\sigma$ = std.\ dev.}
    \label{tab:word2vec-results}
\end{table}

When PCA dimensionality reduction was done, it seems SVM performed better than the other models in F1 and Macro-F1, 
and just barely loosing out to Random Forest in Accuracy. I believe the dimensionality reduction allowed the 
tree based methods to perform better as there was less information to deal with and they could split on 
certain features with greater information gain, which is likely why they won in with t-SNE over SVM. 

Overall, PCA has higher values and wins because the dimensionality reduction is not as severe as t-SNE's. 
PCA preserves a sort of global ``structure'' which helps with classification, but t-SNE is more for 
visualization and twists the distances of ``meanings'' at the global/vocabulary level.

And, a final note, that these word-based embeddings \textbf{all} perform worse that frequency based embeddings. 
This could be for a few reasons, given that this method should do better than task 1. 

\begin{itemize}
    \item \textbf{Tweets (or Xeets if you must)}: Given that tweets are generally shorter and don't contain as much context as 
    other documents do, and that Word2Vec was potentially trained on more cohesive and well structured documents,
    this could have an effect.
    \item \textbf{Irony}: The entire dataset is designed around detecting irony, which has subverted 
    expectations in the texts, causing the ``nonsequitors'' of irony to trip up the models
    \item \textbf{Preprocessing}: This is most unlikely, but there's a chance that either I did 
    not do enough or did too much preprocessing on the data, removing or keeping information 
    that ultimately hurt the models after the embeddings were done.
\end{itemize}

\section{Sentence Embeddings} \label{sentence_embed}

For the sentence embeddings, I used the pre-trained Sentence-BERT (SBERT) model. I used Sentence-BERT because we 
had talked about BERT in class, and I was more familiar with the idea overall beyond the 
Universal Sentence Encoder (USE) model.

For preprocessing, I used only 2 of the options I had listed earlier. SBERT was trained on 
natural prose and it relies on word order, punctuation, and un-lemmatized words to gain more information.
So, the 2 preprocessing steps that I used were (3) Remove URLs \& Mentions and (4) Strip Hashtag Symbol. This 
was minimal and done only to affect the nature and structure tweets (Xeets, if you must).

As we did not talk about BERT or SBERT in class yet, but it uses concepts that we talked about in 
class (like the window of looking at 2 words to the left and right of the target word), I will do 
my best to give a brief overview of SBERT.

\begin{center}
    BERT reads text in both directions, not just left to right. This lets BERT capture a lot more 
    context that earlier models \cite{devlin-etal-2019-bert}. But, BERT was sadly not designed to do sentence level embeddings. 
    So, in 2019 SBERT was released to use siamese and triplet network structures on sentence \textit{pairs} 
    \cite{reimers-gurevych-2019-sentence}. This allows two semantically similar sentences 
    to be represented close to each other (geometrically speaking).
\end{center}

\begin{figure}[H]
    \centering
    \includegraphics[width=\textwidth]{images/results_sbert.png}
    \caption{Sentence-BERT embedding results.}
    \label{fig:sbert}
\end{figure}

\begin{table}[H]
    \centering
    \renewcommand{\arraystretch}{1.4}
    \begin{tabular}{l
                    >{\centering\arraybackslash}p{1.8cm}
                    >{\centering\arraybackslash}p{1.8cm}
                    >{\centering\arraybackslash}p{1.8cm}
                    >{\centering\arraybackslash}p{1.8cm}
                    >{\centering\arraybackslash}p{1.8cm}
                    >{\centering\arraybackslash}p{1.8cm}}
        \toprule
        \rowcolor{tableheader}
        \color{white}\textbf{Classifier}
            & \color{white}\textbf{Acc.}
            & \color{white}\textbf{Acc. $\sigma$}
            & \color{white}\textbf{F1}
            & \color{white}\textbf{F1 $\sigma$}
            & \color{white}\textbf{Macro F1}
            & \color{white}\textbf{MF1 $\sigma$} \\
        \midrule
        \rowcolor{rowalt} Random Forest  & 0.6407 & 0.0069 & 0.5691 & 0.0114 & 0.6305 & 0.0078 \\
                          SVM            & \cellcolor{green!13}\textbf{0.6772} & 0.0094 & \cellcolor{green!13}\textbf{0.6484} & 0.0098 & \cellcolor{green!13}\textbf{0.6750} & 0.0093 \\
        \rowcolor{rowalt} Decision Tree  & 0.5499 & 0.0133 & 0.5404 & 0.0137 & 0.5496 & 0.0133 \\
        \bottomrule
    \end{tabular}
    \caption{Sentence-BERT embedding results (5-fold CV). Acc.\ = Accuracy, MF1 = Macro F1, $\sigma$ = std.\ dev.}
    \label{tab:sbert-results}
\end{table}

Well, SVM is now back on top! 

Decision Tree crumbled under pressure here, doing very poorly with the dense vectors that SBERT produces and 
not having the ensemble like Random Forest to help it potentially intialize better. So, SVM crushes them 
because, like stated in section \ref{freq_embed}, SVM does better in high dimensional problems. 

Also, I believe that all of these values (for SVM) are better than in both section \ref{freq_embed} and 
\ref{word_embed} because SBERT produces ``higher quality'' vectors that contain lots of information 
between embeddings whereas the frequency and word based embeddings were missing out on a lot of 
larger, document wide context. This means that the same word, like ``bank'' from the in class examples, 
has different meaning based on its surrounding words, which SBERT captures, but BOW and TF-IDF can only 
represent one word, one way. It's also worth noting that compared to the word-based embeddings task, SBERT 
does not need dimensionality reduction to perform and outperform, which is very practical in real world 
use cases to not need to perform extra compuations and resources on dim. reduction.

What we should have seen in this report is a steady increasing of results, starting with frequency embedding 
and coming to sentence embedding, with SBERT/USE being on top. But, Word2Vec was an outlier. So, it's 
reassuring to see that SBERT is ultimately better than both of the other approaches.

\section*{Topic Modeling}

Out of interest in the BERT lineage, I decided to go with BERTopic for topic modeling. 

For the topic modeling task, I did everything in one go. I tested the BERTopic model with topics (K) 
ranging from 1 to 10. However, BERTopic showed the behavior of returning one fewer topic than was input because
of BERTopic uses HDBSCAN clustering which could merge topics that are too similar. From my Intro to 
Urban Computing class, as well as Machine Learning 1, DBSCAN clusters points that are similar and if 
certain points are too close together that were topically different, they could then be clustered into 
one topic. So, for our optimal topic number of 6, it only returned 5 different topics. 

\begin{figure}[H]
    \centering
    \includegraphics[width=\textwidth]{images/results_topic_modeling.png}
    \caption{Topic Modeling Results: Optimal number of topics after testing with k=2 - k=10}
    \label{fig:topic_mod_k}
\end{figure}

\begin{table}[H]
    \centering
    \renewcommand{\arraystretch}{1.4}
    \begin{tabular}{>{\centering\arraybackslash}p{2.2cm}
                    >{\centering\arraybackslash}p{2.2cm}
                    >{\centering\arraybackslash}p{4.0cm}}
        \toprule
        \rowcolor{tableheader}
        \color{white}\textbf{$k$} & \color{white}\textbf{Actual Topics} & \color{white}\textbf{Coherence Score} \\
        \midrule
        \rowcolor{rowalt} 2  & 1 & N/A \\
                          3  & 2 & 0.4243 \\
        \rowcolor{rowalt} 4  & 3 & 0.3833 \\
                          5  & 4 & 0.3951 \\
        \rowcolor{rowalt} \textbf{6}  & \textbf{5} & \cellcolor{green!13}\textbf{0.4722} \\
                          7  & 6 & 0.4045 \\
        \rowcolor{rowalt} 8  & 7 & 0.4311 \\
                          9  & 8 & 0.3672 \\
        \rowcolor{rowalt} 10 & 9 & 0.3690 \\
        \bottomrule
    \end{tabular}
    \caption{BERTopic coherence scores for $k = 2$ to $k = 10$ requested topics. Highest coherence at $k=6$ (5 actual topics).}
    \label{tab:topic-modeling}
\end{table}

\begin{figure}[H]
    \centering
    \includegraphics[width=\textwidth]{images/results_topic_modeling_wordclouds.png}
    \caption{Topic Modeling Results: Word Cloud produced from BERTopic}
    \label{fig:topic_mod_word}
\end{figure}

From the results, the coherence score is best when topics $k = 6$ which returns 5 topics in total. Then, 
from those 5 returned topics, we can generate our word clouds for each topic. 

The word clouds that are produced seem to have somewhat recognizable topics, but with a less than .5 
coherence score it is difficult to discern clear topics. Topic 3 is politically and racially motivated 
but does not have a clear grounding whereas topic 5 has a more prominent grounding of "notices within the eu".
Topic 1 has a difficult to discern central topic from all the rest, seeming to either be general irony and 
sarcasm but it is unclear. Finally, topic 2 is revolving around teams and games, and topic 4 is revolving 
around media and consuming media like music and television. 

\section*{Findings}

Across all three embedding approaches, a winner can be crowned! Sentence-BERT (SBERT) with SVM achieved 
the highest results overall (Acc: 0.6772, F1: 0.6484, Marco F1: 0.6750), followed by the frequency based 
methods, and finally Word2Vec performing worst of the three methods. 

The gap between BOW and TF-IDF was small but still noticable, and I believe that is because stop word 
removal during preprocessing already eliminated a lot of common words that TF-IDF would remove, so I 
inadvertently helped BOW be more competitive with TF-IDF.

Word2Vec performing worse than BOW and TF-IDF is an outlier. The most likely causes were: the short 
and informal nature of tweets conflicting with the Wikipedia-trained model's vocabulary and the 
difficulty irony poses for word level embeddings because ironic language deliberately subverts 
the expected outcome, and in this case semantic relationships the Word2Vec relies on and was likely 
trained on. PCA consistently outperformed t-SNE, as t-SNE's 2 dimensional compression discards too much 
global context.

For topic modeling, BERTopic with an optimal $k = 6$ after sweeping through 9 different $k$ values 
(which only yielded 5 actual topics after HDBSCAN) had the best coherence of 0.4722. The five topics 
were rather broad and have a good amount of room for interpretation, but generally revolved around irony 
itself, politics, race, media consumption, and the EU. When the best coherence we find is below 0.5, it 
shows that our dataset proved difficult in extracting clean topics from short, informal, and intentionally 
misleading/subversive text like tweets (Xeets, if you must).

\section*{Code}

This code is linked \href{https://drive.google.com/file/d/1EARU1KkUAKrpfGEqAsisFIZriJhQWj2b/view?usp=sharing}{here}.

\bibliography{sources}

\end{document}